\documentclass{article} % For LaTeX2e
\usepackage{iclr2027_conference,times}

\input{math_commands.tex}

\usepackage{hyperref}
\usepackage{url}
\usepackage{booktabs}
\usepackage{amsmath}
\usepackage{amssymb}
\usepackage{graphicx}

\title{Adversarial Skill Tuning: Self-Play Between\\ Natural-Language Policies for Mathematical Error Detection}

\author{Anonymous}

%\iclrfinalcopy

\begin{document}

\maketitle

\begin{abstract}
Adversarial self-play is an attractive recipe for training verifiers: a perturber corrupts
correct reasoning, a verifier hunts for the corruption, and neither needs human error
annotation. In practice the recipe is brittle --- the game is difficult to balance, and
reward hacking is endemic. We ask whether the same game can be played at the level of
\emph{natural-language policies} rather than model weights. We instantiate a perturber and
a verifier as \emph{agent skills}: versioned prose documents that a coding agent reads and
follows. After each round of self-play the agent inspects the scored episodes and rewrites
the losing policy, producing \texttt{perturb-v\{N+1\}} and \texttt{verify-v\{N+1\}}, in
direct analogy to a policy-gradient step but with the update expressed in text. Rewards are
computed by the same programmatic reward function an RL implementation would use, so the
two regimes optimise an identical objective. We separate each skill into an
\emph{invariant} contract that the updater may not touch and a \emph{tuned} policy segment
that it rewrites, and we require the updater to file reward-earning behaviours that exploit
the scorer as bugs rather than lessons. Across three rounds on Hendrycks MATH, self-play
metrics move as expected, and we evaluate every verifier version on held-out ProcessBench,
which the tuning loop is structurally prevented from seeing. Transfer is real but not
monotone: the round-2 policy improves the ProcessBench F1 from $0.869$ to $0.897$, while
the round-3 policy regresses to $0.874$ --- an over-fit to the current perturber that the
self-play reward alone reported as a gain. We argue that a mandatory external benchmark is
not an optional add-on to self-play but the only signal that distinguishes learning to
verify from learning one's opponent.
\end{abstract}

\section{Introduction}

Mathematical and statistical errors survive peer review at a rate that has motivated a
small industry of automated checkers. Recent reporting describes volunteer and commercial
projects applying large language models to published literature at scale, uncovering
arithmetic and methodological mistakes that human referees missed, while cautioning that
the tools themselves produce unverified flags~\citep{gibney2025errors}. The underlying
capability --- reading a piece of reasoning and locating the step where it first goes wrong
--- is also the capability that scalable oversight of language models
requires~\citep{amodei2016concrete,bowman2022measuring,saunders2022self}, and it is
measured directly by ProcessBench~\citep{zheng2024processbench}.

Training such a verifier is bottlenecked by supervision: annotated process errors are
expensive, and existing corpora are small relative to the diversity of ways a derivation can
fail. Adversarial self-play offers an escape. A \emph{perturber} takes a correct solution
and injects errors, producing ground truth for free; a \emph{verifier} tries to recover
them; the perturber is rewarded when the verifier fails and vice versa. The structure is
that of a generative adversarial network~\citep{goodfellow2014generative} transposed onto
discrete text, and of self-play in games~\citep{silver2017mastering}.

Our own experience implementing this loop with reinforcement learning --- GRPO
\citep{shao2024deepseekmath} for the perturber, DPO \citep{rafailov2023direct} for the
verifier --- was that the machinery works and the \emph{game} does not. Across three
training rounds the perturber found three successive reward hacks: declaring errors it had
not injected, injecting a single mistake and declaring it $k$ times, and finally producing
text that was redundant rather than wrong. Each was patched in the reward function, at the
cost of a round of GPU time to discover it. This is the standard reward-hacking
pathology~\citep{skalse2022defining}: the agent optimises the measurement, not the task.

This paper asks whether the same adversarial game becomes easier to steer when the policies
are \emph{prose}. We implement the perturber and verifier as \emph{agent skills} --- markdown
documents that a coding agent \citep{anthropic2025claudecode} loads and follows as its
policy for the round. After each round, a static \emph{updater} skill reads the scored
episodes and rewrites the policy text. The artefact being optimised is human-readable at
every step: one can see the policy that produced a reward, and see the edit that a round of
evidence motivated. Our contributions are:

\begin{itemize}
\item \textbf{A formalisation of adversarial skill tuning} (\S\ref{sec:method}) as
alternating coordinate ascent over a pair of natural-language policies, with the reward
computed by exactly the programmatic function an RL implementation would use.
\item \textbf{Three structural guardrails} that make the loop auditable: an
invariant/tuned split within each skill, mechanically enforced; a
\emph{faithfulness filter} requiring the updater to file scorer-exploiting wins as bugs
rather than convert them to policy; and enforced context isolation between the two players.
\item \textbf{An external-validity protocol} on ProcessBench that the tuning loop cannot
read, together with initial results over three rounds (\S\ref{sec:results}) showing that
self-play gains transfer in round 2 and fail to transfer in round 3.
\end{itemize}

\section{Related work}

\paragraph{Process supervision and error identification.}
Outcome-supervised verifiers~\citep{cobbe2021training} were superseded by process
supervision, where each reasoning step carries a label~\citep{lightman2023verify}, and by
process reward models trained with automatically synthesised step
labels~\citep{wang2024mathshepherd}. ProcessBench~\citep{zheng2024processbench} evaluates
the resulting critics directly: given a problem and a step-segmented solution, a model must
return the index of the earliest erroneous step or declare the chain correct. Its 3,400
human-annotated cases span four difficulty tiers built from
GSM8K~\citep{cobbe2021training}, MATH~\citep{hendrycks2021measuring},
OlympiadBench~\citep{he2024olympiadbench} and Omni-MATH~\citep{gao2024omnimath}. We adopt
both its task definition and its metric.

\paragraph{Automated error detection in scientific writing.}
The applied motivation for verifier research is documented by \citet{gibney2025errors},
who describes projects screening the published literature with language models --- one
approaching authors privately after analysing several hundred papers, another publicly
flagging candidate errors across tens of thousands --- and the attendant risk of unconfirmed
flags. That risk is precisely a precision problem, and it is why our reward and our
benchmark metric both penalise false alarms as heavily as misses.

\paragraph{Self-play and self-improvement.}
Self-play removes the annotation bottleneck by having a model generate its own
curriculum~\citep{silver2017mastering,chen2024selfplay}, and self-rewarding
schemes~\citep{yuan2024selfrewarding} and self-refinement~\citep{madaan2023selfrefine} let a
model supply its own feedback. Our setting differs in that the reward is not model-judged:
it is a deterministic program over the perturber's declared ground truth, which is what
makes the perturber's incentive to cheat concrete and detectable.

\paragraph{Optimising prompts and skills instead of weights.}
A line of work treats the prompt as the parameter, optimising it by search
\citep{zhou2023large}, by natural-language ``gradients'' \citep{pryzant2023automatic}, by
using the model as an optimiser over its own instructions \citep{yang2024large}, or by
compiling declarative pipelines \citep{khattab2024dspy}. Closest to us is
Ctx2Skill \citep{si2026context}, which extracts natural-language \emph{skills} from long
contexts through multi-agent self-play: a Challenger poses probing tasks, a Reasoner
answers, a Judge scores, and Proposer/Generator agents synthesise skill updates from
failures, with a \emph{Cross-Time Replay} mechanism that re-evaluates historical skill
candidates against curated probe sets to prevent adversarial collapse.

We share the premise --- that a natural-language skill is a trainable object, and that
adversarial interaction can produce one without human annotation --- and differ in three
ways. First, our reward is a programmatic function of a machine-checkable ground truth
rather than a model judge, so credit assignment is exact and the failure mode we must guard
against is scorer exploitation rather than judge bias. Second, both of our players are
skills and both are updated, so the curriculum and the solver co-evolve. Third, where
Ctx2Skill guards against collapse by replaying historical candidates on curated in-domain
probes, we guard against it with an external human-annotated benchmark that the update
mechanism is structurally forbidden to read; as \S\ref{sec:results} shows, that distinction
is load-bearing, because the round in which our loop over-fits is invisible from inside the
game.

\section{Method}
\label{sec:method}

\subsection{The adversarial episode}

Let $\mathcal{D}$ be a corpus of solved problems, and let $(p, s) \sim \mathcal{D}$ be a
problem with a correct step-by-step solution. Fix a budget $k \in \mathbb{N}$ of errors per
episode. A \emph{policy} in our setting is a natural-language document $\sigma$; writing
$\pi(\cdot \mid \cdot\,; \sigma)$ for the behaviour induced when an agent follows $\sigma$,
one episode is

\begin{align}
(\hat{s}, E) &\sim \pi_P\!\left(\cdot \mid p, s, k\,;\, \sigma_P\right),
&& E = \{e_i\}_{i=1}^{k},\quad e_i = (o_i, \iota_i, \tau_i, \rho_i), \label{eq:perturber}\\
C &\sim \pi_V\!\left(\cdot \mid p, \hat{s}\,;\, \sigma_V\right),
&& C = \{c_j\}_{j=1}^{m},\quad c_j = (q_j, x_j, \varsigma_j). \label{eq:verifier}
\end{align}

The perturber returns a corrupted solution $\hat{s}$ together with a machine-readable log
$E$ in which each error names the original span $o_i$, the injected span $\iota_i$, a type
$\tau_i$ drawn from a fixed 26-way taxonomy of student misconceptions
\citep{rittlejohnson2025detecting}, and a rationale $\rho_i$. The verifier sees only $p$
and $\hat{s}$ --- never $E$ --- and returns claims $c_j$ consisting of a quoted span $q_j$,
an explanation $x_j$ and an optional type $\varsigma_j$.

\subsection{Matching, error units, and reward}

Scoring requires aligning claims to injected errors. Let $\mathrm{IoU}(\cdot,\cdot)$ be
character-level intersection-over-union after normalisation (LaTeX escaping, math
delimiters, alignment characters, whitespace). Define the alignment score
\begin{equation}
\mu(e, c) \;=\; \max\Big\{
\underbrace{\mathrm{IoU}(\iota_e, q_c)}_{\text{span overlap}},\;
\underbrace{\gamma \cdot \mathbb{1}\!\left[\,q_c \succ \delta(o_e,\iota_e)\,\right]}_{\text{diff coverage}},\;
\underbrace{\lambda(\iota_e, q_c)}_{\text{containment}}
\Big\},
\label{eq:mu}
\end{equation}
where $\delta(o_e,\iota_e)$ is the set of fragments a character diff marks as changed,
$q_c \succ \delta$ holds when the quote is anchored inside the error span \emph{and}
contains a changed fragment, $\gamma = 0.9$, and $\lambda$ is the length ratio of the
shorter normalised string when one contains the other. The middle term matters: it
distinguishes a claim that quotes the actual mistake from one that quotes an untouched
clause of the same sentence.

A greedy one-to-one assignment $M \subseteq E \times C$ then takes, for each error in turn,
the highest-scoring unused claim above a threshold $\theta = 0.5$.

\paragraph{Error units.} A perturber can cap the verifier's recall at $1/k$ by injecting one
mistake and declaring it $k$ times --- as a root error plus its propagated consequences, as
overlapping rewrites of one region, or as pure restatement. We therefore quotient $E$ before
computing recall. Let $\sim$ be the reflexive-symmetric-transitive closure of the relation
that holds between two errors when they occupy overlapping or containing spans, share a
changed diff fragment, differ by a token-Jaccard below $0.6$, or when one modifies a
$\backslash\texttt{boxed}\{\}$ answer downstream of the other. The \emph{error units} are the
quotient $\mathcal{U} = E/\!\sim$, and we restrict to units actually present in $\hat{s}$:
\begin{equation}
\mathcal{U}^{\mathrm{eff}} = \{\, U \in \mathcal{U} \;:\; \exists\, e \in U,\; \iota_e \sqsubseteq \hat{s} \,\},
\qquad
k_{\mathrm{eff}} = \big|\{ e \in E : \iota_e \sqsubseteq \hat{s}\}\big|,
\end{equation}
with $\sqsubseteq$ denoting normalised containment. Recall is measured over units and
precision over claims:
\begin{equation}
\mathrm{rec} \;=\; \frac{\big|\{\,U \in \mathcal{U}^{\mathrm{eff}} : \exists\, e \in U,\ (e,\cdot) \in M \,\}\big|}{\big|\mathcal{U}^{\mathrm{eff}}\big|},
\qquad
\mathrm{prec} \;=\; \frac{|M|}{|C|}.
\label{eq:recprec}
\end{equation}
Quotienting only the recall denominator is deliberate: a verifier that flags both a root
error and its propagation keeps full precision, while a perturber that stacks gains nothing.

\paragraph{Rewards.} With $F_\beta$ the usual harmonic combination of Eqs.~(\ref{eq:recprec}),
\begin{align}
r_V &= F_\beta(\mathrm{prec}, \mathrm{rec})
\;-\; \lambda_{\mathrm{spam}} \max\!\big(0,\ |C| - \varrho\, \max(1,k_{\mathrm{eff}})\big)
\;-\; \lambda_{\mathrm{rep}} \mathbb{1}[\text{output collapsed}], \label{eq:rv}\\[2pt]
r_P &= \big(1 - \mathrm{rec}\big)
\;-\; \lambda_{\mathrm{miss}}\big(k - k_{\mathrm{eff}}\big)
\;-\; \lambda_{\mathrm{dup}}\, d(E)
\;-\; \lambda_{\mathrm{hist}}\, d\big(E, E^{<t}\big), \label{eq:rp}
\end{align}
where $d(E)$ counts near-verbatim repeats within the episode and $d(E, E^{<t})$ counts
injections that are $\geq 0.85$ similar to one used in an earlier round. Both are overridden
by a dominating format penalty when the perturber's output is not a valid instance of the
schema:
\begin{equation}
r_P \;=\;
\begin{cases}
\phi_{\mathrm{json}} = -10, & \text{output is not parseable JSON},\\
\phi_{\mathrm{schema}} = -5, & \text{parses but violates the contract},\\
\text{Eq.~(\ref{eq:rp})}, & \text{otherwise.}
\end{cases}
\label{eq:format}
\end{equation}
The two-level penalty exists because a uniform penalty makes every failing episode score
identically, which in the RL implementation eliminated the within-group reward variance that
GRPO learns from; it plays the same role here by giving the updater a gradient across
failure severities.

\subsection{Skill tuning as a textual policy update}

Each policy document is partitioned into two segments,
\begin{equation}
\sigma \;=\; \big(\, \underbrace{\iota}_{\text{invariant}} \;,\; \underbrace{\varpi}_{\text{tuned}} \,\big),
\end{equation}
where $\iota$ transcribes the I/O contract, the taxonomy and the reward definition from the
implementation, and $\varpi$ holds the strategy. A round $t$ produces a scored dataset
\begin{equation}
\mathcal{E}^t \;=\; \Big\{\big(p_n, s_n, \hat{s}_n, E_n, C_n, M_n, r^P_n, r^V_n\big)\Big\}_{n=1}^{N}.
\end{equation}
The update operator $\Lambda$ is a language model that reads $\mathcal{E}^t$ and the current
policy and emits a new one:
\begin{equation}
\varpi^{t+1} \;=\; \Lambda\big(\varpi^{t},\, \mathcal{A}(\mathcal{E}^t)\big),
\qquad
\iota^{t+1} = \iota^{t},
\label{eq:update}
\end{equation}
subject to $|\mathrm{edits}| \le 5$, each edit citing at least one episode index. Equation~(\ref{eq:update}) is the textual analogue of a policy-gradient step: $\mathcal{E}^t$ plays
the role of the sampled batch, the updater's diagnosis plays the role of credit assignment,
and the edit budget plays the role of a trust region. Alternating the two updates,
\begin{equation}
\varpi_P^{t+1} = \Lambda_P\big(\varpi_P^{t}, \mathcal{A}(\mathcal{E}^t)\big),
\qquad
\varpi_V^{t+1} = \Lambda_V\big(\varpi_V^{t}, \mathcal{A}(\mathcal{E}^t)\big),
\end{equation}
is coordinate ascent on the pair of objectives
$J_P(\sigma_P,\sigma_V) = \mathbb{E}[r_P]$ and $J_V(\sigma_P,\sigma_V) = \mathbb{E}[r_V]$
of the underlying min-max game. A \texttt{freeze} flag disables either updater, recovering
the ablation in which one player is held fixed.

\paragraph{The faithfulness filter.} The operator $\mathcal{A}$ in Eq.~(\ref{eq:update}) is the
crux. Let $X^t \subseteq \mathcal{E}^t$ be the episodes whose reward was earned by
exploiting the scorer rather than the task --- for the perturber, stacking, restatement, or
burying a small change inside a long declared span; for the verifier, quoting broadly to
maximise overlap, shotgunning claims, or keying on the opponent's habits. Then
\begin{equation}
\mathcal{A}(\mathcal{E}^t) \;=\; \mathcal{E}^t \setminus X^t,
\end{equation}
and $X^t$ is written to a report as a set of candidate defects in the reward implementation.
Without $\mathcal{A}$, skill tuning would rediscover in prose exactly the hacks that RL
found in weights, and would do so faster, because a language model is better at reading a
reward function than a gradient is. The test the updater is instructed to apply is not
``did this score well'' but ``would a competent mathematician call this a real mistake, and
a hard one''.

\paragraph{Enforced isolation.} Because a single agent plays both roles, leakage is the
default failure. The harness makes it structural: after the perturber writes its output, a
program validates it and emits a file containing only $(p, \hat{s})$; the verifier pass runs
in a fresh agent context scoped to a directory holding nothing else. A round in which the
isolation cannot be established is discarded rather than scored.

\paragraph{Convergence.} The loop halts when the round budget is exhausted, or when every
active updater reports that $\mathcal{A}(\mathcal{E}^t)$ contains no admissible edit, or
when the round-over-round change in the relevant metric stays inside a band
($|\Delta r_P| < 0.05$, $|\Delta F_1^{\text{self-play}}| < 0.05$) for two consecutive
rounds.

\section{Experimental setup}
\label{sec:setup}

\paragraph{Implementation.} Both players and both updaters are skills in the Claude Code
agent harness \citep{anthropic2025claudecode}, executed by Claude Opus 5. Episode sampling,
perturbation validation, reward computation and aggregation are a Python harness that shares
its reward module with the repository's RL trainers, so a skill-tuning episode and an RL
episode are scored by identical code.

\paragraph{Self-play rounds.} Each round draws $N = 8$ problems from the Hendrycks MATH
training split \citep{hendrycks2021measuring} with $k = 3$ errors per episode, sampling
without reuse across rounds. Neither player is frozen. We report three completed rounds; a
fourth is in progress and is not analysed here.

\paragraph{External evaluation.} After each round the newly produced verifier is evaluated
on ProcessBench \citep{zheng2024processbench}. We sample $20$ items per subset, balanced
between erroneous and correct chains, fixed by a seed that does not depend on the round, so
all versions are scored on the same $80$ held-out items. Each item is presented as a single
text with explicit boundaries,
\begin{equation*}
\texttt{<step\_0>}\, s_0 \,\texttt{</step\_0>} \;\cdots\; \texttt{<step\_{n-1}>}\, s_{n-1} \,\texttt{</step\_{n-1}>},
\end{equation*}
so the whole chain is judged in context while every claim remains addressable to a step. The
verifier's claim set is reduced to a ProcessBench prediction by
\begin{equation}
\widehat{y} \;=\;
\begin{cases}
-1, & C = \emptyset,\\
\min_{c \in C} \mathrm{step}(c), & \text{otherwise},
\end{cases}
\label{eq:pred}
\end{equation}
where $\mathrm{step}(c)$ is the claim's declared index, or the index of the block containing
its quoted span when the declaration is missing or out of range. A claim set that resolves
to no step is scored wrong against both classes: accusing something unlocatable is not a
detection. Following \citet{zheng2024processbench} we report the accuracies on erroneous and
correct chains and their harmonic mean,
\begin{equation}
\mathrm{Acc}_{\mathrm{err}} = \frac{1}{|\mathcal{I}_{-}|}\sum_{i \in \mathcal{I}_{-}} \mathbb{1}[\widehat{y}_i = y_i],
\quad
\mathrm{Acc}_{\mathrm{cor}} = \frac{1}{|\mathcal{I}_{+}|}\sum_{i \in \mathcal{I}_{+}} \mathbb{1}[\widehat{y}_i = -1],
\quad
F_1 = \frac{2\,\mathrm{Acc}_{\mathrm{err}}\mathrm{Acc}_{\mathrm{cor}}}{\mathrm{Acc}_{\mathrm{err}} + \mathrm{Acc}_{\mathrm{cor}}}.
\label{eq:pbf1}
\end{equation}
The harmonic mean is essential rather than cosmetic: a verifier that never claims anything
attains $\mathrm{Acc}_{\mathrm{cor}} = 1$ and a balanced accuracy of $0.5$, but $F_1 = 0$.

\paragraph{Leakage control.} ProcessBench is a test set for this loop and is never read by
it. Items and gold labels live in separate directories; the evaluation runs in a fresh agent
scoped to the items; the summary artefact contains identifiers, hit/miss flags and
aggregates but no item text and no labels; and both updaters are explicitly forbidden from
reading the evaluation tree. No ProcessBench content enters a skill, a rollout log, or the
orchestrating context at any point.

\section{Results}
\label{sec:results}

\subsection{Self-play dynamics}

Table~\ref{tab:selfplay} reports the three completed rounds. Format validity is perfect
throughout --- the contract segment of the perturber skill states the schema constraints
explicitly, and unlike the RL perturber, which failed the schema in $71\%$ of round-1
episodes, the skill perturber never violates it. The adversarial pressure is visible in the
trend: mean $r_P$ rises from $0.125$ to $0.313$ while verifier recall falls from $0.875$ to
$0.688$, i.e.\ the perturber is winning ground over the three rounds.

\begin{table}[t]
\caption{Self-play rounds. Each round is 8 episodes with $k = 3$ on Hendrycks MATH; both
players update after every round. Recall is over error units, Eq.~(\ref{eq:recprec}).
FP counts verifier claims matching no injected error.}
\label{tab:selfplay}
\begin{center}
\small
\begin{tabular}{lccccccccc}
\toprule
Round & Perturber & Verifier & Fmt.\ valid & $r_P$ & $r_V$ & Recall & Prec. & Units/ep & FP \\
\midrule
1 & \texttt{v1} & \texttt{v1} & 8/8 & 0.125 & 0.900 & 0.875 & 0.938 & 2.75 & 1 \\
2 & \texttt{v2} & \texttt{v2} & 8/8 & 0.167 & 0.833 & 0.833 & 0.833 & 3.00 & 4 \\
3 & \texttt{v3} & \texttt{v3} & 8/8 & 0.313 & 0.750 & 0.688 & 0.875 & 2.88 & 2 \\
\bottomrule
\end{tabular}
\end{center}
\end{table}

The qualitative record is more informative than the aggregates. Round 1 produced five
perturber edits and five verifier edits, all traceable to specific episodes: the verifier's
anti-propagation rule had been over-applied, skipping two independently invalid lines
because they \emph{looked} like consequences of an earlier error, and was rewritten to
demand an explicit re-derivation test. Round 2's perturber update shifted strategy from
arithmetic to \emph{definitional} errors --- corrupting which quantity a symbol denotes or
what a ratio compares --- on the observation that these leave every later line internally
consistent, so re-derivation alone cannot surface them. Round 3's verifier update responded
by promoting a ``model audit against the problem statement'' to equal standing with
re-derivation. That edit is the one that fails to transfer.

\subsection{Transfer to ProcessBench}

Table~\ref{tab:processbench} reports each verifier version on the same 80 held-out items.

\begin{table}[t]
\caption{ProcessBench, all three tuned verifier versions, $n = 80$ (20 per subset, balanced
erroneous/correct, identical items across versions). $F_1$ is Eq.~(\ref{eq:pbf1}); ``Pooled''
computes it over all 80 items, ``Macro'' averages the four subset $F_1$ values as in
\citet{zheng2024processbench}.}
\label{tab:processbench}
\begin{center}
\begin{tabular}{lcccccc}
\toprule
Verifier & GSM8K & MATH & Olympiad. & Omni-MATH & Macro & Pooled \\
\midrule
\texttt{verify-v2} (after R1) & 0.889 & 0.889 & 0.889 & 0.800 & 0.867 & 0.869 \\
\texttt{verify-v3} (after R2) & 0.889 & \textbf{0.947} & \textbf{0.947} & 0.800 & \textbf{0.896} & \textbf{0.897} \\
\texttt{verify-v4} (after R3) & 0.889 & \textbf{0.947} & 0.847 & 0.800 & 0.871 & 0.874 \\
\bottomrule
\end{tabular}
\end{center}
\end{table}

Round 2's edits transfer cleanly: pooled $F_1$ rises $0.869 \rightarrow 0.897$. Round 3's do
not: $0.897 \rightarrow 0.874$. At the item level the movement is unusually clean in both
directions. From \texttt{v2} to \texttt{v3}, exactly two items change and both change from
wrong to right (one MATH, one OlympiadBench), with no regressions. From \texttt{v3} to
\texttt{v4}, exactly two items change and both change from right to wrong, both on
OlympiadBench, with no gains.

\subsection{Anatomy of the round-3 regression}

The loss is confined to a single cell. Error detection does not move between \texttt{v3} and
\texttt{v4} ($\mathrm{Acc}_{\mathrm{err}} = 0.85$ for both); OlympiadBench
$\mathrm{Acc}_{\mathrm{cor}}$ falls from $1.00$ to $0.80$. The v4 policy kept every detection
v3 had and added two false alarms on clean chains.

This is the predictable consequence of the two v3$\rightarrow$v4 edits, which pushed in the
same direction: one gave the definitional audit equal weight with re-derivation, the other
asked for more separately-anchored claims. The verifier's own round report describes the
mechanism without knowing it was being measured on it --- it flagged chains that ``reach a
correct final answer through an invalid step\ldots\ following the rule that a chain can be
arithmetically fine and still be wrong about the model''. Against the perturber, that rule
earned recall: it was written from an episode in which two genuine definitional errors went
unclaimed. Against human-annotated chains, ProcessBench's labels disagree with enough of
those judgements to cost more than the rule gains.

This is the outcome the external evaluation exists to detect, and the self-play reward could
not have detected it: from inside the game, round 3 looked like a verifier that had closed a
real gap. Note the shape of the failure. It is not that skill tuning does not work --- round
2's edits produced a genuine $+0.029$ --- but that the round in which the perturber shifted
strategy is the round in which the verifier over-fit to that strategy.

\subsection{Context: published ProcessBench baselines}

Table~\ref{tab:baselines} reproduces reference points from \citet{zheng2024processbench}.
These are \emph{not} directly comparable to Table~\ref{tab:processbench} and we do not claim
a state-of-the-art result: the published figures use the full 3,400-case benchmark with its
natural class balance, our figures use 80 balanced items; the baselines are 2024-era models
prompted as single-pass critics, while our verifier is a 2026 frontier model following a
tuned multi-step policy inside an agent harness. The table establishes only the scale of the
task and the difficulty gradient across subsets.

\begin{table}[t]
\caption{Published ProcessBench $F_1$ (\%) from \citet{zheng2024processbench}, Table 3, on
the full benchmark. Shown for scale only; see the caveats in the text.}
\label{tab:baselines}
\begin{center}
\small
\begin{tabular}{lccccc}
\toprule
Model & GSM8K & MATH & Olympiad. & Omni-MATH & Average \\
\midrule
Qwen2.5-Math-7B-PRM800K (PRM) & 68.2 & 62.6 & 50.7 & 44.3 & 56.5 \\
GPT-4o-0806 (critic) & 79.2 & 63.6 & 51.4 & 53.5 & 61.9 \\
QwQ-32B-Preview (critic) & 88.0 & 78.7 & 57.8 & 61.3 & 71.5 \\
o1-mini (critic) & 93.2 & 88.9 & 87.2 & 82.4 & 87.9 \\
\bottomrule
\end{tabular}
\end{center}
\end{table}

Two observations survive the caveats. First, the difficulty gradient reported for every
model in \citet{zheng2024processbench} --- monotone decline from GSM8K to Omni-MATH --- is
only weakly present in our runs, where Omni-MATH is the weakest subset ($0.800$ throughout)
but MATH and OlympiadBench are not separated from GSM8K. Second, Omni-MATH is the one subset
that no round moved at all, which is consistent with its errors being of a kind the
perturber never produces.

\subsection{Reward-implementation defects surfaced by the filter}

The faithfulness filter $\mathcal{A}$ logged 18 candidate scorer defects across the six
update reports of three rounds --- about ten distinct issues, since the two updaters
independently flag the same matcher artifacts --- none of which entered a policy. The dominant class is span-alignment brittleness: in three
round-2 episodes a correct detection failed to match its ground truth by a one- or
two-character margin (a trailing period, a closing bracket) that normalisation does not
strip, and in one round-3 episode a claim that correctly diagnosed a false identity scored
$0.441$ and $0.467$ against the two errors it spanned --- both just under the $\theta = 0.5$
gate. The aggregate effect is material: round 2's recorded $-0.067$ verifier regression is,
on the updater's corrected reading, approximately $+0.058$. We report the uncorrected numbers
in Table~\ref{tab:selfplay} and treat the discrepancy as a measurement limitation rather than
adjusting the results. The methodological point is that the filter converts what would have
been silent reward-hacking pressure into an auditable list of bugs.

\section{Limitations}

The evaluation is small. Eighty items, ten per class per subset, means each subset accuracy
moves in increments of $0.1$ and each reported transition is a two-item change; Wilson
intervals on the pooled accuracies span roughly $\pm 0.12$, so neither the $+0.029$ nor the
$-0.023$ is individually significant. We report them because the item-level direction is
unanimous in both cases, not because the sample supports a statistical claim. Scaling to the
full 3,400 cases is the obvious next step.

We have no measurement of the untuned seed policy. The first evaluation was run after round
1, so \texttt{verify-v1} has no ProcessBench score, and we cannot separate the quality of the
hand-written seed from the contribution of tuning. Self-play rounds are likewise small (8
episodes), and only one seed and one problem source were used.

Finally, the perturber's error distribution is not the distribution ProcessBench samples
from. Injected misconceptions in Hendrycks MATH solutions are, by construction, localised and
of known type; ProcessBench errors arise naturally in model-generated chains and include
omissions and invalid framings that our taxonomy does not generate. Transfer across that gap
is precisely what we measure, and the round-3 result shows the gap is wide enough to punish
a policy that specialises.

\section{Conclusion}

We formulated adversarial self-play over natural-language policies, in which a perturber and
a verifier are versioned skills read by an agent, scored by the same programmatic reward an
RL implementation would use, and rewritten between rounds by a static updater. Three
structural choices --- an invariant contract the updater cannot edit, a filter that converts
scorer-exploiting wins into bug reports, and enforced context isolation between the players
--- make the loop auditable in a way that a weight-space implementation of the same game is
not: every policy, every edit, and every rejected lesson is a readable artefact tied to the
episodes that motivated it.

Our initial results are two rounds of transfer and one round of over-fitting. We regard the
second as the more useful finding. Self-play produced a verifier that was measurably better
at beating its own opponent and measurably worse at the task, and no quantity available
inside the game distinguished the two. Held-out human-annotated evaluation is therefore not
a reporting convenience for adversarial skill tuning; it is the only available definition of
success.

\bibliography{arappav}
\bibliographystyle{iclr2027_conference}

\appendix
\section{Appendix: the edit that did not transfer}
\label{app:diff}

The verifier policy is a numbered list of rules in the tuned segment $\varpi_V$. Below is
rule \textbf{V3} as it stood in \texttt{verify-v3} (ProcessBench $F_1 = 0.897$) and after the
round-3 update produced \texttt{verify-v4} ($F_1 = 0.874$). The update was motivated by an
episode in which two genuine definitional errors went unclaimed; it is reproduced verbatim
except for line wrapping.

\paragraph{\texttt{verify-v3}.}
\begin{quote}\small
\textbf{V3 --- Sweep for the standard misconceptions, then check the model itself.} Sign
handling across a move between sides; numerator/denominator or dividend/divisor order; common
denominators; additive reasoning where the relation is multiplicative; off-by-one in sequence
indices; probability outside $[0,1]$; a step that quietly stops short of what the problem
asked. Add the ones that hide best: an inverse operation applied in the wrong direction, and
an identity or factorisation whose \emph{form} is invalid while its constants look plausible.
Then do a pass that arithmetic cannot do for you --- read every definition against the
problem statement. \ldots\ An error planted in a definition leaves every later line
consistent with it, so re-derivation returns nothing; it is invisible except by comparison
with the problem.
\end{quote}

\paragraph{\texttt{verify-v4}.}
\begin{quote}\small
\textbf{V3 --- Sweep the misconceptions, then audit the model against the problem statement.}
\ldots\ Then do the pass that arithmetic cannot do for you, \textbf{and give it the same
weight as the re-derivation, because it is where the errors you miss actually live}. Go back
to the problem statement and check the solution's model against it line by line: what each
symbol denotes, what a named ratio compares, which set or range a stated condition actually
picks out, which quantity is being minimised or maximised, and whether the thing finally
reported is the thing that was asked for. A wrong \emph{set membership} or a wrong
\emph{objective} leaves every subsequent computation internally flawless, so no amount of
re-derivation will surface it --- a solution can be arithmetically perfect and still answer
the wrong question. \ldots
\end{quote}

The promotion from ``also check the model'' to ``give it the same weight'' is the whole
regression. Against the round-3 perturber, which had itself just specialised in definitional
corruption, the rule recovered exactly the errors the previous version missed. On
ProcessBench it produced two false alarms on OlympiadBench chains that are arithmetically
sound and whose framing the annotators accepted, and no compensating detections: error
accuracy is unchanged at $0.85$, correct-chain accuracy on that subset falls from $1.00$ to
$0.80$.

Two properties of this failure are worth noting for the design of such loops. First, it is
\emph{legible}: the responsible edit is a single paragraph in a versioned text file, cited to
the episode that motivated it, which is not true of the corresponding weight update in an RL
implementation. Second, it is \emph{invisible from inside the game} --- the same edit raised
self-play recall on the round that motivated it, and the faithfulness filter did not catch it
because the behaviour is not scorer exploitation at all. It is a correct generalisation from
the opponent's distribution to a policy, and the opponent's distribution was simply not the
world's.

\end{document}
